% Reptation sliding-length theory — insert into main PRL manuscript
% Replace $g$ with $\sqrt{A^*}$ or $\Omega^*$ as appropriate.

\subsection*{Reptation sliding length}

\paragraph{Physical picture.}
When a rod is confined within a cylindrical tube, it undergoes a series of
grazing collisions with the tube wall.  Each collision is characterised by an
elastic normal rebound (coefficient of restitution $e=1$) and Coulomb wall
friction: the tangential (axial) velocity component is reduced by
$\Delta v_{\parallel} = 2\mu v_{\perp}$ per collision, where $v_{\perp}$ is
the instantaneous speed normal to the wall and $\mu$ is the kinetic friction
coefficient.  Denoting the radial gap between the rod surface and the tube wall
as $g$, and the rod's initial axial and transverse speeds as $v_{\parallel}$
and $v_{\perp}$, we compute the total axial displacement $\langle L \rangle$
the rod travels before its axial momentum is exhausted.

\paragraph{Discrete theory.}
The rod starts at the tube axis and traverses an initial one-way leg of
duration $g/v_{\perp}$, covering axial distance $v_{\parallel}g/v_{\perp}$,
before its first wall contact.  After each subsequent bounce $k$
($k=1,2,\ldots$), the axial speed becomes
$v_{\parallel}^{(k)} = v_{\parallel} - 2k\mu v_{\perp}$, and the rod
traverses the full tube diameter in time $2g/v_{\perp}$.  Summing all
productive legs gives
%
\begin{equation}
  \langle L \rangle
  = \frac{g}{v_{\perp}}
    \!\left[
      v_{\parallel}
      + 2\sum_{k=1}^{N}\!\left(v_{\parallel} - 2k\mu v_{\perp}\right)
    \right]
  = \frac{g}{v_{\perp}}
    \!\left[
      v_{\parallel}(1+2N) - 2\mu v_{\perp} N(N+1)
    \right],
  \label{eq:reptation_L}
\end{equation}
%
where $N = \lfloor v_{\parallel}/(2\mu v_{\perp})\rfloor$ is the number of
bounces and the arithmetic series has been evaluated explicitly.

\paragraph{Two regimes.}
Two distinct regimes are separated by a critical friction coefficient
$\mu_c = v_\parallel/(2v_\perp)$.  For $\mu \geq \mu_c$, a single bounce
drains all axial momentum ($N=0$), giving
%
\begin{equation}
  \langle L \rangle = \frac{v_\parallel}{v_\perp}\,g
  \qquad (\mu \geq \mu_c),
  \label{eq:reptation_single}
\end{equation}
%
independent of $\mu$.  In the many-bounce limit $\mu \ll \mu_c$,
Eq.~\eqref{eq:reptation_L} reduces to an arithmetic-series average,
%
\begin{equation}
  \langle L \rangle \approx \frac{v_\parallel^2}{2\mu v_\perp^2}\,g
  \qquad (\mu \ll \mu_c).
  \label{eq:reptation_multi}
\end{equation}
%
This is a factor of two smaller than the naive estimate obtained by ignoring
the velocity decrement at each bounce.  For an isotropic initialisation
($v_\parallel \sim v_\perp$), this asymptote becomes
$\langle L \rangle \approx g/(2\mu)$.

Later, $g$ will be replaced by $\sqrt{A^*}$ (the effective confinement radius
derived from the mean tube cross-sectional area $A^*$) or $\Omega^*$ (a
geometry-weighted tube diameter), connecting this single-rod result to the
collective reptation statistics of a dense rod assembly.

\paragraph{Simulation comparison.}
We validate Eq.~\eqref{eq:reptation_L} with rigid-rod simulations using a
non-smooth contact (NSC) solver with elastic rebounds and Coulomb friction.  A
single rod of length $\ell=1$ (aspect ratio $\ell/d=100$) is launched along
the axis of a cylinder over a grid of gaps
$g \in [0.001,\,0.10]$ and friction coefficients $\mu \in [0.11,\,1.0]$
(90 parameter pairs), with equal initial axial and transverse speeds
($v_\parallel = v_\perp$, giving $\mu_c = 0.5$).
Figure~\ref{fig:reptation_theory} shows the normalised sliding length
$L/\ell$ plotted against $g/(\mu\ell)$.  In the multi-bounce regime
($\mu < \mu_c$, circles), all data collapse onto the universal line
$L/\ell = \tfrac{1}{2}\,g/(\mu\ell)$; in the single-bounce regime
($\mu \geq \mu_c$, squares) the sliding length saturates at $L/\ell = g/\ell$
(horizontal dashed lines).  The full discrete prediction
Eq.~\eqref{eq:reptation_L} reproduces all 90 simulation results to within
$0.6\%$ r.m.s.\ and $3.9\%$ maximum error.

\begin{figure}[htbp]
  \centering
  \includegraphics[width=0.55\textwidth]{results/reptation_theory_collapse}
  \caption{%
    Normalised axial sliding length $L/\ell$ as a function of $g/(\mu\ell)$
    from rigid-rod simulations (symbols), coloured by $\mu$.
    Circles: multi-bounce regime ($\mu < \mu_c = 0.5$);
    squares: single-bounce regime ($\mu \geq \mu_c$).
    Solid line: asymptotic theory $L/\ell = \tfrac{1}{2}\,g/(\mu\ell)$,
    Eq.~\eqref{eq:reptation_multi}.
    Dashed horizontal lines: saturation values $g/\ell$ for each gap tested,
    Eq.~\eqref{eq:reptation_single}.
  }
  \label{fig:reptation_theory}
\end{figure}
